\documentclass{article}
\usepackage[frenchb]{babel}  % Nécessaire pour les exposants dans du texte
\usepackage[utf8]{inputenc}
\usepackage[justification=centering]{caption} % Pour centrer le texte des légendes qui sont sur plusieurs lignes
\usepackage{amsmath}  % Fonctionnalités mathématiques
\usepackage{amsfonts,amssymb}  % Typographie des symboles mathématiques
\usepackage{multirow}   % Fusion de cellules de tableau
\usepackage{subfigure}   % Alignement des images

\DeclareMathOperator{\sinc}{sin_c}  % Fonction sinus cardinal


\title{Module hyperfréquence et optoélectronique \\ \bigskip \bigskip \hrule width 16.5cm \bigskip \textbf{Rapport de TP \\ Analyse vectorielle de dispositifs passifs \\} \bigskip \hrule width 16.5cm}
\author{Gabriel-Robez Baptiste}
\date{20/01/2020}

\pdfinfo{
  /Title    (Module hyperfréquence et optoélectronique - Rapport de TP)
  /Author   (Gabriel-Robez Baptiste)
  /Subject  ()
  /Keywords (hyperfréquence, ADS)
}
\usepackage[utf8]{inputenc}
\usepackage{graphicx}
\usepackage[top=2cm, bottom=3cm, left=2.5cm, right=2.5cm]{geometry}

\begin{document}
\maketitle

\renewcommand{\contentsname}{Sommaire}  % pour remplacer "Table des matières" par "Sommaire"
\tableofcontents

\bigskip

\begin{figure}[!ht]
\centering
\includegraphics[scale=0.15]{logo_phelma.png}
\end{figure}

\clearpage




\section{Introduction}

\section{Conclusion}


\end{document}
